% !TeX root = ../教材题库_全部条目_转写工作稿.tex
% Source: XJ-B1-C04-EX-P102-Y0669
% 第4章 / 4.1 实数指数幂和幂函数 / 教材 p.102、103 / PDF 第109、110页
% OCR未校对：公式与文字均可能有误；未完成数学TeX转写；禁止直接发给学生
% 原题及相邻原解截图（核对用）：../题目原图/XJ-B1-C04-EX-P102-Y0669-01.webp, ../题目原图/XJ-B1-C04-EX-P102-Y0669-02.webp
% 本文件只在转写完成后作为题干使用；OCR 条目尚不是可靠题干。
\begin{bankitem}{4.1 例题 例题9}{第4章 / 4.1 实数指数幂和幂函数 / 教材 p.102、103 / PDF 第109、110页}{XJ-B1-C04-EX-P102-Y0669}{OCR未校对：公式与文字均可能有误；未完成数学TeX转写；禁止直接发给学生}
\begin{ocrdraft}
[教材p.102 / PDF第109页]
\par
0若丆0）是幂函数，
\par
且丆（8）一16，求丿、（9），不64L
\par
解设丆0）一．
\par
由已知条件得不8）一8“
\par
=16=8\ensuremath{\times}2=8\ensuremath{\times}83=8
\par
4
\par
故
\par
，所以不工）一产
\par
3
\par
于是丆（9）一9六不64）一643一
\par
=28=256．
\par
下面举例说明幂函数在实际生活中的应用．
\par
大雾天，海陆空的交通运输都会受到影响．
\par
我们知道，雾是大量小水滴在空气中悬浮而形成
\par
的．小水滴为什么不掉下来呢？下面来分析．
\par
把水滴近似看成小球，设半径为则它的
\par
体积V=—mx\ensuremath{\cdot}
\par
在物理学中，由于重力G=mg，m=pV'于
\par
是G=pgV,其中p是水的密度（常数），g=9.8m/s2.于是水滴所受的重力G和水
\par
滴的体积V成正比，又V一一3，于是重力G与成正比．
\par
另一方面，水滴所受的空气阻力丆与水滴的表面积s成正比，又s一4“2，因
\par
而阻力丿、与成正比．
\par
于是，当半径从1mm变小到0．1mm时，水滴所受的重力G减小到原来的
\par
1
\par
，而阻力丆减小到原来的丽．相对来说，阻力与重力的比值增加到10倍．同
\par
1000
\par
理，工从1mm变小到0．01mm时，阻力与重力的比值增加到原来的100倍．因此，
\par
[教材p.103 / PDF第110页]
\par
随着半径的变小，水滴所受的重力G比起阻力丆来，很快就变得微不足道了．
\par
这
\par
就是小水滴不掉下来形成雾的原因\ensuremath{\cdot}
\par
又如，由物理学知识可知，在用点光源照明时，照
\par
度E与被照亮的面到光源的距离R的平方成反比，即
\par
E一（I为发光强度).
\par
又如，根据航天技术的最基本的理论依据一一一一一万有
\par
引力定律可知，与地心距离为7．的物体所受
\par
的引力的大小与严成反比．
\par
再如，悬浮的颗粒在不流动的空气中扩散的速度。
\par
与时间的算术平方根孬成正比，这是一种缓慢的运动．这就说明了为什么工业生产
\par
过程中排放的烟尘，在没有风的情况下，会长时间笼罩在一片地方的上空．
\end{ocrdraft}
\end{bankitem}
